% ============================================================
%  discussion.tex  —  §8
%  Target venue: Omega
%  Updated: 29/Sep/2026
% ============================================================

\section{Discussion}
\label{sec:discussion}

\subsection{Two mechanisms, ordered by severity}

The results separate two distinct ways in which a normalisation choice
degrades portfolio quality, and the distinction matters because the
remedies differ.

The first is \emph{implicit re-weighting}, characterised in closed form
by Lemma~\ref{lem:amplitude}: an additive score built on normalised
columns of unequal amplitude operates on effective weights
$w^{\mathrm{eff}}\neq w$, whatever the analyst elicited.
The effect is graded, is summarised by the total
variation~\eqref{eq:TV_results}, and orders the three normalisations
monotonically within additive aggregation (Table~\ref{tab:tvorder}).
It is also invisible: nothing in a standard TOPSIS or SAW report reveals
that \texttt{perm} entered at $2.87$ times its nominal weight.

The second is \emph{score--evaluator mismatch}.
Proposition~\ref{prop:p1} guarantees Pareto efficiency for exact
maximisation of an additive score only when the score and the value
functions defining the objectives use the same normalisation.
When they do not, the guarantee is not weakened but voided, and the
observed cost exceeds what re-weighting alone predicts by roughly a
factor of seven (Section~\ref{sec:refset}).
Two findings are attributable to this mechanism: the collapse under
reference-set misalignment (Table~\ref{tab:fourthcell}), and the
near-absence of any gain from exact assembly under vector normalisation
($+2.2$~pp, interval containing zero).

The practical ordering follows.
Re-weighting is a modelling error that can be quantified and, if
accepted, documented.
Mismatch is a specification error that removes the theoretical warrant
for optimising at all.

\subsection{Relation to portfolio decision analysis}

The closest antecedents study portfolio construction from
multi-criteria scores.
\citet{BransMareschal1992} introduced constrained selection over
PROMETHEE flows; \citet{Mavrotas2006} combined multi-criteria scores
with binary programming; \citet{VetscheraDeAlmeida2012} formulated
$c$-optimal portfolios, of which the present problem with $\kk=3$ is an
instance; and \citet{DeAlmeidaVetschera2012} showed that PROMETHEE~V
portfolios depend on the scale transformation applied to the flows.

The last of these is the nearest precedent, and the difference is
instructive.
\citet{DeAlmeidaVetschera2012} establish that scale matters for one
method's portfolios, without an external reference against which to
measure the consequence.
We give a closed form for the induced re-weighting, applicable to any
additive aggregation and any normalisation, and measure the consequence
against an exhaustively enumerated Pareto set.
The robust portfolio modelling line \citep{Liesio2007,Liesio2008} and
SMAA-PO \citep{Yang2015} compute non-dominated portfolios exactly, but
in a scalar value space under weight uncertainty; our reference set
retains three unaggregated partial value functions, which permits the
dominance test to detect trade-off losses that a scalar criterion cannot.

Neither line imposes a hard pairwise separation constraint.
That constraint is standard in the location literature
\citep{MoonChaudhry1984,Erkut1996}, and it is what makes the assembly
effect combinatorial rather than parametric: a site foreclosed by greedy
selection cannot be recovered by any re-weighting
(Section~\ref{sec:greedy}).

\subsection{Relation to the normalisation and rank-reversal literature}

Set dependence of normalisation is documented, principally for TOPSIS.
\citet{GarciaCascales2012} demonstrate rank reversal under changes to
the alternative set and recommend scales that do not depend on it;
\citet{Kuo2017} identifies a failure of the closeness index;
\citet{OpricovicTzeng2004} observe that VIKOR uses a linear and TOPSIS
a vector normalisation, with consequences for unit invariance.
Empirical comparisons proceed largely by rank-agreement statistics
\citep{Jahan2015,Aytekin2021,Vafaei2018,Vafaei2022,Salabun2020Sym}.

Our contribution is to give the mechanism a closed form and to measure
its cost against an exact reference rather than against another ranking.
Three previously separate observations become corollaries of one
identity: the vector-normalisation distortion of
Proposition~\ref{prop:p3}, the tie-induced contraction in rank space
(Remark~\ref{rem:p2normalization}), and the amplitude loss from
normalising over inadmissible alternatives~\eqref{eq:refsetA}.
The account also explains a structural fact the rank-agreement literature
records without explanation: VIKOR's invariance to vector normalisation
follows immediately, since a common amplitude factor cancels
in~\eqref{eq:weff}.
\todo{Before submission: targeted search on ``effective weights'',
  ``implicit weighting'', ``normalisation weight distortion'' in MCDM to
  confirm the closed form is not already published. If it is, reposition
  the contribution as the empirical cost measurement against an exact
  frontier, which stands independently.}

\subsection{The PROMETHEE--Borda identity}

Proposition~\ref{prop:p2} is a corollary of results established three
decades ago.
\citet{Bouyssou1992} characterises the net flow method axiomatically and
shows that it reduces to Copeland or Borda when the valued preference
relations are crisp; \citet{Marchant1996} extends the Borda equivalence
to valued relations.
With the usual preference function the relations are crisp by
construction, so the identity is implied.
We restate it to fix the affine constants, to record that it holds
exactly under arbitrary tie structure, and chiefly because it isolates
normalisation as the sole source of empirical divergence: in our data
the entire departure of the implemented rank score from the net flow is
attributable to rescaling midranks by $N+1$ rather than to the
aggregation.
Without the identity, that divergence would plausibly be attributed to
the method.

\subsection{Implications for practice}

Three recommendations follow, in decreasing order of cost to ignore.

First, normalise over the set of alternatives that can actually be
chosen, and report the set used.
Screening criteria applied after normalisation leave the value scale
calibrated on alternatives that are unavailable, which is the most
damaging configuration we observe.

Second, verify that the score and the value model share a normalisation
before adding exact optimisation.
Integer programming on a mismatched score optimises the mismatch; the
$+29.8$~pp available under min--max falls to $+2.2$~pp under vector
normalisation.

Third, report effective weights alongside elicited weights.
The computation is one line given the normalised matrix, and it exposes
distortions that no rank-agreement diagnostic reveals.
Where elicited weights carry normative content --- as with AHP judgements
from stakeholders --- a normalisation that silently reweights them is
difficult to defend.

Comparing aggregation operators, by contrast, is the least consequential
of the design choices we examine ($+0.9$~pp, interval containing zero),
and the mechanism explains why: aggregation acts downstream of the
distortion.

\subsection{Limitations}

\emph{Scope of the closed form.}
Lemma~\ref{lem:amplitude} is exact for additive aggregation.
TOPSIS and VIKOR aggregate distances to ideal points, where amplitudes
enter quadratically or in ratio form; the mechanism persists but the
identity does not.
Extending the closed form to distance-based aggregations is the natural
next step.

\emph{Empirical support for the ordering.}
The monotone relation between predicted distortion and observed
performance rests on three normalisation levels within SAW.
Establishing a predictive coefficient requires five or more levels, which
additional normalisation families (sum, max, Weitendorf) would supply at
modest computational cost.

\emph{Single instance.}
One candidate set, one weight vector source, $\kk=3$.
The analytic results are instance-independent, but the magnitudes depend
on the marginal distributions; \texttt{perm} with 350 tied zeros is an
unusually adversarial case.
Fogo Island quantifies the mechanism; it does not estimate its typical
size.

\emph{Weight perturbation model.}
Dirichlet perturbation around $\wahp$ represents symmetric uncertainty in
relative weights.
For structured disagreement among stakeholders,
SMAA-style weight-space analysis \citep{Lahdelma2001} would be
appropriate.

\emph{Regulatory grounding.}
The exclusion of the $\Nexcl$ overlapping sites is treated as a
screening decision.
\todo{D2: with the Cape Verde instrument cited, this becomes a normative
  constraint and the reference-set argument strengthens accordingly.}
